% P2 methods upgrade: independent publication claim and no patient identity claim.
\subsection{Validation units, overlap graph and three evaluation estimands}
\label{sec:p2-units}
A cross-person EEG result can involve four non-equivalent units:
a clinical individual, a released digital participant ID,
a continuous raw recording, and processed labelled EEG
epochs. Only digital recording contents can be compared
without further information from the clinical data stewards.
It follows that exact cross-ID numerical identity does
\emph{not} establish identity of two patients.

Under fixed preprocessing let $h(x,y)$ denote an exact
fingerprint of a trial's numeric array and event label.
Released identifiers $g$ and $g'$ are joined if their
processed trial sets share a fingerprint. The connected
components define \emph{recording-content groups}.
If $g$ is held out by its released ID but $g'$ in the
training set shares all of $g$'s processed labelled EEG
trials, a nominal participant-independent evaluation
is not independent at the level of digital signal
contents. A content-disjoint split instead removes
every member of the held content component from
training. The graph construction is deterministic and
standard; its scientific value is in making the
claimed independence unit inspectable.

We distinguish three reporting estimands:
\[
\begin{split}
 M_{50}&=\frac{1}{50}\sum_{g\in\mathrm{released IDs}}A_g,
 \\
 M_{49,A}&=\frac{1}{49}\sum_{G\in\mathrm{content groups}}
                  A_G(F_A,\lambda_A),\\
 M_{49,B}&=\frac{1}{49}\sum_{G\in\mathrm{content groups}}
                  A_G(F_B,\lambda_B).
\end{split}
\]
$M_{50}$ uses nominal ID weights, whereas $M_{49,A}$
uses the original trained models and original source
choices on 49 unique evaluated contents. Thus
$M_{49,A}-M_{50}$ is an evaluation \emph{weighting}
contrast. $M_{49,B}-M_{49,A}$ is a matched-target
\emph{software pipeline} contrast with changes to
source fitting and parameter selection.
Their sum is arithmetic, not an identified
causal patient-level leakage effect.

\subsection{Distinguishing our question from previous leakage research}
\label{sec:p2-related-novelty}
Brookshire et al.\ \cite{brookshire2024leakage}
showed that segment-wise rather than participant-wise
validation can dramatically overestimate unseen-person
EEG classification. Gil and Hern\'andez-Sabat\'e
\cite{gil2026subjectvalidation} compared multiple
levels of EEG validation across six classifiers.
Del Pup et al.\ \cite{delpup2025partition} compared five EEG
validation strategies, three classification problems and
four deep architectures in a much broader study of
partition choices. These studies demonstrate that
validation-unit errors are not merely an issue of one EEG algorithm; however, they did not test exact cross-ID signal-content aliasing as the unit definition itself.
MOABB \cite{jayaram2018moabb} promotes a transparent
and repeatable public-EEG evaluation interface.
Our study asks the narrower additional question of
whether the released \emph{participant identifier}
itself correctly keys independent digital EEG contents
before such a split is constructed.
Neither the cryptographic fingerprint, graph component
procedure nor the factorial subtraction is claimed as
a previously unknown mathematical invention.

\subsection{Negative controls, preprocessing dependence and raw signal check}
\label{sec:p2-controls}
An exact fingerprint depends on signal sampling,
event alignment, filtering, channel order, data dtype,
and the specified MOABB processing path.
Here those choices are pinned, and the total number
of analysed trials is reported explicitly.
We additionally compared the raw public EDF data:
32 of 33 reader-interpreted channel arrays were equal
over the corresponding record pair, whereas one
channel reported an empty name and different values.
This is independent supporting \emph{digital signal}
evidence but not proof of clinical identity or an
explanation of how these files were created.
An independent second raw-reading implementation and
clarification from data custodians are still pending.

Weibo2014 and BNCI2014\_002 were retained as
two previously exposed control acquisitions.
No cross-ID exact processed epoch matches were
found within the indicated pre-processing pipelines.
Zero exact overlaps cannot rule out near duplicates,
window overlaps, or other release errors.
A second negative control arises within Liu2024:
two IDs with coincident historical target AUC
vectors exhibited \emph{different} signal contents.
Equal classifier outcomes are therefore not an
admissible identity fingerprint.

\subsection{Matched model-by-selection design and dependence}
\label{sec:p2-analysis-units}
All 49 \emph{same} held digital recording contents
were evaluated under
$(F_A,\lambda_A)$, $(F_A,\lambda_B)$,
$(F_B,\lambda_A)$ and $(F_B,\lambda_B)$.
We repeat the same comparison for source-macro
discriminative selection, proper source-trained-mean
predictive Gaussian selection and a historical
class-centred covariance score. Each policy's source
parameters were frozen before the crossed held-target
responses were computed. For group $g$ and policy $m$,
\[
 I_{gm}=Y_{BB,gm}-Y_{BA,gm}-Y_{AB,gm}+Y_{AA,gm}.
\]
The fitted source model includes the training-derived
Riemannian tangent geometry; this can change under
content purging even when $\lambda$ does not change.
The full result was recomputed from the original
49-row immutable archive and checked against its
independently frozen paired model scores with
$10^{-10}$ tolerance. Its content SHA-256 was
\texttt{0e7bfdb73970fcff9231a51954bd2d548f472cb296624967afc80d9d241ce2be}.

The 49 held content groups are valid reporting
\emph{units}, not 49 independent trained classifiers:
their source pools overlap almost completely.
No ordinary paired population test on the 49 fitted
AUC differences is therefore claimed here.
Instead we report conditional distributions of
paired changes, source-choice changes and
delete-one-content-group sensitivity.
All original held outcomes had already been
exposed when these descriptive checks were designed;
there is no independent prospective replication.
